% ID: FIS-OND-MEC-001
% Titolo: Misurare l'energia trasportata da un'onda
% Disciplina: Fisica
% Area: Onde
% Argomento: Onde meccaniche
% Sottoargomento: Energia trasportata da un'onda
% Classe: Quarta liceo scientifico
% Difficolta: 4
% Tipo: modellizzazione
% Risultato: misura di lavoro o potenza tramite un galleggiante
% Tempo_stimato: 10 min
% Competenze: progettazione di un esperimento, modellizzazione, uso di schemi ed equazioni
% Prerequisiti: onde meccaniche, energia, ampiezza, potenza
% Tag: fisica, onde, onde_meccaniche, energia, modellizzazione
% Fonte: verifica_fisica_onde_anonimizzata
% Autore: Riccardo Carli
% Licenza: CC-BY-SA-4.0

\begin{esercizio}
Immagina di dover misurare l'energia trasportata da un'onda che si propaga
sulla superficie di un lago.

Descrivi un possibile esperimento adatto allo scopo, facendo uso di schemi ed
equazioni a supporto della tua idea.
\end{esercizio}

\begin{soluzione}
Un possibile esperimento consiste nell'usare un piccolo galleggiante collegato
a una guida verticale e a un sensore di forza, oppure a una molla di costante
elastica nota \(k\). Il passaggio dell'onda fa oscillare verticalmente il
galleggiante.

Se si usa una molla, misurando lo spostamento massimo \(A\) del galleggiante si
puo stimare l'energia elastica massima accumulata:
\[
E_\text{max}=\frac{1}{2}kA^2.
\]
Questa energia proviene dall'onda che ha compiuto lavoro sul galleggiante.

Un metodo piu generale e misurare nel tempo la forza \(F(t)\) esercitata
dall'onda e lo spostamento verticale \(y(t)\) del galleggiante. Il lavoro
scambiato in un ciclo e:
\[
L=\oint F\,dy.
\]
Se il periodo dell'onda e \(T\), la potenza media trasferita al dispositivo e:
\[
P_\text{media}=\frac{L}{T}.
\]

Ripetendo la misura e conoscendo la larghezza efficace \(b\) del dispositivo,
si puo stimare anche l'energia trasportata per unita di tempo e per unita di
fronte d'onda:
\[
\frac{P_\text{media}}{b}.
\]
L'esperimento richiede di ridurre gli attriti e di calibrare il sensore, perche
una parte dell'energia viene dissipata.
\end{soluzione}
